% !TEX root = ../linnik399.tex
\section{The near rows}\label{sec:rows}

A \emph{near row} is a constraint obtained by applying the threshold inequality
\eqref{eq:T} to selected characters and zeros. In a leaf program (\cref{subsec:vocabulary}),
characters are grouped by parameter into bins, each with its own column. Known entries from the
first or reserved family may instead be recorded separately as \emph{family terms}. There are
four kinds of rows, the family, shifted, graded and two-test rows; \cref{tab:rowkinds} in
\cref{subsec:rowkinds} lists their anchors, sieve weights and entries.

We first specify the test functions and normalization. We then justify each permitted
kind of entry and describe how to combine them into rows. Most rows follow from
\cref{cor:zeroform}; the two-test row in \cref{subsec:twotest} has a separate proof.
\Cref{prop:rowsvalid} is the conclusion needed later: every actual zero configuration
in a leaf satisfies each of its near rows at some threshold.

\subsection{The parabolic test function}
All detectors and Gram tests are normalized parabolic autocorrelations. For $c>0$ put
\[
  P(u)=(1-u)^3(1+3u+u^2)=1-5u^2+5u^3-u^5,\qquad
  \hat f_c(t)=\begin{cases}P(t/c),&0\le t\le c,\\0,&t\ge c.\end{cases}
\]
Thus $\hat f_{2\gamma}=f_\gamma/f_\gamma(0)$ with $f_\gamma$ as in \eqref{eq:ftest}. Write $\hat F_c$ for its
transform.

\begin{lemma}[Properties of the parabolic test functions]\label{lem:parabolic}
Let $c>0$.
\begin{enumerate}[(i)]
  \item $\hat f_c\ge0$, $\hat f_c(0)=1$, $\hat f_c$ is decreasing on $[0,c]$, and $\hat f_c$ extended by
    $0$ is twice continuously differentiable on $[0,\infty)$ with $|\hat f_c''|\le10/c^2$. So
    Condition~1 holds.
  \item $\Re\hat F_c(iy)=60c\,(cy\cos\frac{cy}2-2\sin\frac{cy}2)^2/(cy)^6\ge0$ for $y\ne0$, and
    Condition~2 holds.
  \item $x\mapsto\hat F_c(x)$ is positive and strictly decreasing on $\RR$.
  \item For $z\ne0$, $\hat F_c(z)=c\,\Phi(cz)$, where
    \[
      \Phi(w)=\frac1w-\frac{10}{w^3}+\frac{30(1+e^{-w})}{w^4}
        +\frac{120e^{-w}}{w^5}-\frac{120(1-e^{-w})}{w^6}\qquad(w\ne0).
    \]
    The singularity at $0$ is removable. Consequently, for $d\ge0$ and $|y|\ge U>0$,
    \[
      \Re\hat F_c(-d+iy)\le cR(cU,e^{cd}),\qquad-\Re\hat F_c(-d+iy)\le\frac d{U^2}+cR(cU,e^{cd}),
    \]
    where $R(u,E)=\frac{10}{u^3}+\frac{30}{u^4}+\frac{120}{u^6}+E\bigl(\frac{30}{u^4}+\frac{120}{u^5}+\frac{120}{u^6}\bigr)$.
  \item For $d\ge0$ put $\mathcal C_c(d)=\sup_{\Re z\ge-d}\bigl(-\Re\hat F_c(z)\bigr)$. It is finite,
    \[
      \mathcal C_c(d)=\max\Bigl(0,\ \sup_{y\in\RR}\bigl(-\Re\hat F_c(-d+iy)\bigr)\Bigr),
    \]
    and $\mathcal C_c(0)=0$.
\end{enumerate}
\end{lemma}

\begin{proof}
(i) is elementary: $P'(u)=-5u(2+u)(1-u)^2$, and $P(1)=P'(1)=P''(1)=0$. (ii) The function
$\hat f_c$ is a positive multiple of the autocorrelation of $x\mapsto(\gamma^2-x^2)_+$ with $c=2\gamma$, so
its transform on the imaginary axis is a positive multiple of $(\int_0^\infty(\gamma^2-x^2)_+\cos(xy)\,dx)^2$
\cite[Lemma~3.6]{Xylouris2011nullstellen}, which gives the closed form. Condition~2 then follows from
\cref{inp:comparison} with $F_2=0$, since $\hat F_c(z)\to0$ uniformly in $\Re z\ge0$. (iii)
$\hat F_c'(x)=-\int t\hat f_c(t)e^{-xt}\,dt<0$. (iv) Five integrations by parts, using the values of
$P$ and its derivatives at $0$ and $1$. For $z=-d+iy$ we have $c\Re\frac1{cz}=\frac{-d}{d^2+y^2}$, which is
$\le0$ and $\ge-d/U^2$, and each other term is bounded by the corresponding term of $cR(c|z|,e^{cd})$,
which decreases in $|z|\ge U$. (v) $-\Re\hat F_c(z-d)$ is harmonic in $\Re z>0$, continuous up to the
boundary and tends to $0$ uniformly at infinity, so its supremum over $\Re z\ge0$ is attained on
the boundary or is $0$; $\mathcal C_c(0)=0$ by Condition~2.
\end{proof}

\subsection{Row data}\label{subsec:rowdata}
Every family, shifted and graded row fixes rational numbers $\gamma_f>0$ (detector), $\gamma_g>0$ (Gram test),
$s_1\le s$ (safe and reference anchors), $t_0$ and $\mu_{\mathrm{row}}$, and uses the near datum of \cref{subsec:neardata} with
$f=\hat f_{2\gamma_f}$, $g=\hat f_{2\gamma_g}$, $\eps'=\frac1{2000}$, and $2000$ equal cells
$t_0<t_1<\dots<t_{2000}=2\gamma_f$, where $t_0>\frac13+2\eps'$. The heights $h_k\ge0$ are rational numbers.
They are chosen by a heuristic that brings $\omega$ close to the Cauchy--Schwarz optimum
$\omega\propto e^{st}f$: the height of cell $k$ is the positive part of
$e^{sm_k}f(m_k)/\sqrt{\mu_{\mathrm{row}}\kappa_k}-g(m_k)e^{s_1m_k}$ at the midpoint $m_k$, where
$\kappa_k=(t_{k+1}^2-t_k^2)/(2\vartheta_k(t_{k+1}-t_k))$ is the mean sieve cost per unit length of the cell,
computed in binary64 arithmetic and converted to a rational number with denominator at most
$10^{12}$. Any nonnegative heights are admissible, so this choice plays no role in the proofs. In
every family and shifted row $\mu_{\mathrm{row}}=10^9$ and all the computed heights are $0$, so $\Xi\equiv0$. The
offsets are the finitely many numbers $s-\sigma$ that occur for the anchors $\sigma$ of the entries. In
every row the lower bound of $\omega$ on each cell of the support of $f$ is checked to be positive; in
the rows with $\Xi\equiv0$ this holds because $\gamma_g>\gamma_f$.

We use the normalization of \cref{cor:zeroform}: $I_u$ is an upper Riemann
sum for $I_B$ over $4000$ cells, computed with outward interval arithmetic; $D_u\ge D(0)$; each entry
$j$ receives the feature $v_j=r_j/\sqrt{I_uD_u}-\eta'$ (rounded down, and replaced by $0$ if negative)
and the diagonal $(1+\eta')D_j^+/D_u$ (rounded up), where $D_j^+$ is an upper bound for
$D(\delta_j)-d_1+e_j$; and the row has the correlation term $d=(1+\eta')(d_1/D_u+\eta')$ (rounded up). Here
$\eta'=10^{-6}$. For an entry with offset $\delta$ the value $D(\lfloor\delta\rfloor_{1/200})\ge D(\delta)$ is used
(here $\lfloor\delta\rfloor_{1/200}$ is $\delta$ rounded down to a multiple of $\frac1{200}$, and $D$ is decreasing), and an
entry is omitted if its diagonal numerator does not meet the positivity margin used
in the construction. An omitted entry has feature $0$ in the row and contributes
nothing to \eqref{eq:T}. The rounding directions are justified by the last assertion of
\cref{cor:zeroform}.

\subsection{Facts available in a leaf}\label{subsec:leaffacts}
A leaf specifies intervals and lower bounds for the first few zeros. The full definition
of a specification $\mathfrak s$ and its configuration set $C(\mathfrak s)$ is in
\cref{def:specification}; the information needed for this section is listed below.
Write $\nu(\mathcal F)$ for the least parameter in a family among zeros of physical
height at most $1$, and $\nu_{T^*}(\chi)$ for the corresponding minimum at height at
most $T^*$. A reserved family minimizes $\nu$ outside the first family. It need not
be the family attaining the global parameter $\lambda_2$.

Fix a leaf with first-zero cell $[a,b]$ and a configuration in $C(\mathfrak s)$.
For all sufficiently large $q$, the following facts are available.
\begin{enumerate}[(Z1)]
  \item $\lambda_1\in[a,b]$. The \emph{safe anchor} of the leaf is $s_1=\lfloor50a\rfloor/50\le\lambda_1$.
  \item $\lambda'\ge p^*$, where $p^*=\lo'$ if there is a gap (an interval $[\lo',\hi']$ known to
    contain $\lambda'$) and $p^*=p$ otherwise; if there is a finite gap, $\lambda'\le\hi'$.
  \item Every zero in $R(l)$ of a character outside the first family has parameter at least
    $\lambda_2\ge l_2$.
  \item Every family $\mathcal F\ne \mathcal F_1$ has $\nu(\mathcal F)\ge r$, and ordinary characters have $\nu_{T^*}\ge\nu\ge r$.
  \item If $\mathfrak s$ is reserved, the reserved family has $n_2$ characters with common height-one
    parameter $\nu_*\in[\lo_2,\hi_2]$.
  \item Inside: $|\mu_1|\le C_B$. Outside: $|\mu_1|>C_B$.
\end{enumerate}

\begin{lemma}[No zero to the right]\label{lem:noright}
Let $\chi$ be nonprincipal, let $|\gamma|\le2l$ and $0<\delta\le1$, and fix a real anchor
$\sigma$. For sufficiently large $q$, every zero $\rho$ of $L(s,\chi)$ in the disc
$|1+i\gamma-\rho|\le\delta$ has $\lambda_\rho\ge\sigma$ in either of the following cases:
\begin{enumerate}[(a)]
  \item $\sigma\le\lambda_1$;
  \item $\chi\notin\mathcal F_1$ and $\sigma\le l_2$, where $l_2\le\lambda_2$ is the leaf's
    global second-family lower bound.
\end{enumerate}
Equivalently, this disc contains no zero to the right of the line
$\Re s=1-\sigma/\LL$.
\end{lemma}

\begin{proof}
Such a zero has $|\Im\rho|\le2l+\delta<10l$. By \cref{inp:rectangle} it lies in $R(l)$ or has
$\lambda_\rho>\frac13\log\log\LL$, which exceeds every fixed $\sigma$. In $R(l)$ its parameter is at least
$\lambda_1$, and at least $\lambda_2\ge l_2$ if $\chi\notin \mathcal F_1$ (Z3).
\end{proof}

\begin{lemma}[Separation of zeros beyond the representative strip]\label{lem:strip}
Let $\chi$ be an ordinary character with $T^*$-representative $\rho_\chi$, of height $\gamma_\chi$ and
parameter $\lambda_\chi$, and let $\sigma\le\min(\lambda_\chi,s_{\max})$. Then every zero $\rho$ of $L(s,\chi)$ with
$|1+i\gamma_\chi-\rho|\le\delta_0$ and $\lambda_\rho<\sigma$ satisfies $|\gamma_\chi\LL-\mu_\rho|\ge M$, and there are at
most $K_0$ such zeros, counted with multiplicity.
\end{lemma}

\begin{proof}
Such a zero has $|\Im\rho|\le T^*+\delta_0<2\le10l$ and $\lambda_\rho<\sigma\le s_{\max}<\frac13\log\log\LL$, so it
lies in $R(10l)$ and hence, by \cref{inp:rectangle}, in $R(l)$. As $\lambda_\rho<\sigma\le\lambda_\chi$, the minimality
of the representative gives $|\Im\rho|>T^*$; so the zero is counted in $K_0$. By \cref{lem:Tstar},
$|\Im\rho|\ge T^*+M/\LL$, so $\LL|\Im\rho-\gamma_\chi|\ge\LL(|\Im\rho|-T^*)\ge M$.
\end{proof}

\subsection{Entries}\label{subsec:entries}
We now list the entries that occur in the rows, with their kept zeros. In each case the
hypotheses of \cref{lem:response} hold with the stated kept set $\cS$ and anchor $\sigma$; we write
$r$ for the resulting lower bound for the response, before normalization, so that the entry's
feature is $r/\sqrt{I_uD_u}-\eta'$. \Cref{tab:entries} summarizes the entries. Recall that $F(x)$ is
positive and decreasing on $\RR$, and that a kept zero with $\lambda_\rho\ge\sigma$ has
$m_\rho\Re F(\cdot)\ge\Re F(\cdot)\ge0$.

\begin{table}[ht]
\centering
\footnotesize
\setlength{\tabcolsep}{3.5pt}
\renewcommand{\arraystretch}{1.3}
\begin{tabular}{@{}>{\raggedright\arraybackslash}p{3.45cm}>{\centering\arraybackslash}p{1.35cm}
  >{\centering\arraybackslash}p{1.75cm}>{\centering\arraybackslash}p{3.2cm}
  >{\centering\arraybackslash}p{2.35cm}>{\raggedright\arraybackslash}p{1.65cm}@{}}
\toprule
Entry & Anchor $\sigma$ & Kept set $\cS$ & Response bound $r$ & Diagonal numerator & Rows\\
\midrule
(a) ordinary $\chi$, bin $[\lo,\hi)$, at its $T^*$-representative & $\min(\lo,s)$ & $\{\rho_\chi\}$ &
  $F(\hi-\sigma)-\frac16f(0)$ & $D(s-\sigma)-d_1$ & all three\\
(b) inside first family: $\chi_1$ at $\gamma_1$, and $\bchi_1$ at $-\gamma_1$ & $\min(a,s)$ &
  $\{\rho_1\}$, resp.\ $\{\overline{\rho_1}\}$ & $F(b-\sigma)-\frac16f(0)$ & $D(s-\sigma)-d_1$ & family, graded\\
(c) type $\mathrm{rc}$: $\chi_1$ at $\pm\gamma_1$ & $s_1$ & $\{\rho_1,\overline{\rho_1}\}$ &
  $F(b-s_1)+R_{\lo}-\frac16f(0)$ & $D(0)-d_1+(c^{\hi}-d_1)_+$ & family\\
(d) second zero, type $\mathrm{complex}$: $\chi_1$ at $\gamma_1,\gamma'$, $\bchi_1$ at $-\gamma_1,-\gamma'$ &
  $s_1$ & $\{\rho_1,\rho'\}$ and conjugates & $F(b-s_1)+R_p-\frac16f(0)$ at $\gamma_1$;
  $F(\hi'-s_1)+R_q-\frac16f(0)$ at $\gamma'$ & $D(0)-d_1+(c^{\hi}-d_1)_+$ & family\\
(e) shifted first family: $\chi_1$ at $\gamma_1$, and $\bchi_1$ at $-\gamma_1$ & $s$ &
  $\{\rho_1\}$; $\{\rho_1,\overline{\rho_1}\}$ for $\mathrm{rc}$ & $F(b-s)-\1_{\mathrm{rc}}C_Z-\frac16f(0)$ &
  $D(0)-d_1$ & shifted\\
(f) reserved family, at height-one representatives & $\sigma_2$ & the representative &
  $F(\hi_2-\sigma_2)-\frac16f(0)$ & $D(s-\sigma_2)-d_1$ & all three\\
(g) outside first family: $\chi_1$ at $\gamma_1$ (and $\bchi_1$), and one hidden entry per character &
  $s_1$ & $\{\rho_1\}$; $\{\rho_h\}$ & $F(b-s_1)-\frac16f(0)$; $F(\hi_h-s_1)-\frac16f(0)$ & $D(0)-d_1$ &
  family\\
\bottomrule
\end{tabular}
\caption{The entries of the near rows, justified in the list below. The anchor $\sigma$ enters
through the offset $s-\sigma$, where $s$ is the reference anchor of the row; $\sigma_2$ and the
constants $R_{\lo}$, $R_p$, $R_q$, $c^{\hi}$ and $C_Z$ are defined in the list. ``All three'' means
the family, shifted and graded rows of \cref{subsec:rowkinds}.}
\label{tab:entries}
\end{table}

\begin{enumerate}[(a)]
  \item \emph{Ordinary character in the bin $[\lo,\hi)$.} Entry $(\chi,\gamma_\chi,s-\sigma)$ at its
    $T^*$-representative, with $\sigma=\min(\lo,s)$; $\cS=\{\rho_\chi\}$. If $\sigma\le l_2$ there is no
    zero to the right (\cref{lem:noright}); otherwise \cref{lem:strip} gives the hypotheses. Then
    $r=F(\hi-\sigma)-\frac16f(0)$, and the diagonal numerator is $D(s-\sigma)-d_1$.
  \item \emph{Inside first family, one entry per character.} Entries $(\chi_1,\gamma_1,s-\sigma)$ and, for
    type $\mathrm{complex}$, $(\bchi_1,-\gamma_1,s-\sigma)$, with $\sigma=\min(a,s)\le\lambda_1$; $\cS=\{\rho_1\}$
    (resp.\ $\{\overline{\rho_1}\}$). No zero to the right. $r=F(b-\sigma)-\frac16f(0)$, and the diagonal
    numerator is $D(s-\sigma)-d_1$.
  \item \emph{Conjugate pair, type $\mathrm{rc}$} (rows with $s=s_1$ and $\Xi=0$). Entries
    $(\chi_1,\pm\gamma_1,0)$, a pair of entries of the same character with height difference
    $2\gamma_1$; $\cS=\{\rho_1,\overline{\rho_1}\}$ for both. No zero to the right, and
    \[
      r=F(b-s_1)+R_{\lo}-\tfrac16f(0),\qquad R_{\lo}\le\inf\Re F(x+2i\mu)
    \]
    over $x\in[a-s_1,b-s_1]$ and $\mu$ in the height range of the leaf ($[0,1]$ or $[1,\infty)$). Pair
    excess $e\le(c^{\hi}-d_1)_+$ with $c^{\hi}\ge\sup\Re G(-s_1+2i\mu)$ over the range, so the diagonal numerator
    is $D(0)-d_1+(c^{\hi}-d_1)_+$.
  \item \emph{Second zero of the first family} (rows with $s=s_1$ and $\Xi=0$; type $\mathrm{complex}$
    with a finite gap $[p^*,\hi']$). Let $\rho'$ be an admissible zero of $\chi_1$ realizing $\lambda'$
    (\cref{sec:notation}), of height $\gamma'$, namely the one that witnesses the piece of the second-zero split
    (\cref{subsec:subdivision}) in which the configuration lies, and $y=(\gamma_1-\gamma')\LL$, restricted to that
    piece: $|y|\in[0,1],[1,\frac32],[\frac32,2],[2,3],[3,5]$ or $[5,\infty)$.
    Entries $(\chi_1,\gamma_1,0)$, $(\chi_1,\gamma',0)$, $(\bchi_1,-\gamma_1,0)$, $(\bchi_1,-\gamma',0)$, with kept sets
    $\{\rho_1,\rho'\}$, $\{\rho',\rho_1\}$ and their conjugates (a zero is kept only if it lies in the
    disc of the entry). No zero to the right. Then
    \[
      r_{\gamma_1}=F(b-s_1)+R_p-\tfrac16f(0),\qquad r_{\gamma'}=F(\hi'-s_1)+R_q-\tfrac16f(0),
    \]
    where $R_p\le\inf\Re F(x+iy)$ over $x\in[p^*-s_1,\hi'-s_1]$ and $R_q\le\inf\Re F(x+iy)$ over
    $x\in[a-s_1,b-s_1]$, with $|y|$ in the piece. Here $x\ge0$ (as $s_1\le a\le p^*$), so $\Re F(x+iy)\ge0$ by
    Condition~2 and we take $R_p,R_q\ge0$; the multiplicities of the kept zeros can then only increase
    the sums. On the last piece $R_p=R_q=0$, which also covers a zero outside the disc (its term is
    then absent). Each same-character pair is a controlled pair with excess
    $(c^{\hi}-d_1)_+$, $c^{\hi}\ge\sup\Re G(-s_1+iy)$, so every diagonal numerator is $D(0)-d_1+(c^{\hi}-d_1)_+$;
    an entry of $\chi_1$ and an entry of $\bchi_1$ form a pair of distinct characters, with the nonprincipal
    quotient $\chi_1^2$, and carry no mutual excess. If $\rho'=\rho_1$ is a double zero, the two
    entries coincide and $r\ge m_{\rho_1}F(\lambda_1-s_1)-\frac16f(0)$ with $m_{\rho_1}\ge2$, which is at least
    both $r_{\gamma_1}$ and $r_{\gamma'}$ (here $y=0$ and $\lambda'=\lambda_1$).
  \item \emph{Shifted first family} (the shifted row, $s=\min(1.9,p^*,l_2)$, used only if $s>s_1$ and
    $s\ge a$). Entries $(\chi_1,\gamma_1,0)$ and, for type $\mathrm{complex}$, $(\bchi_1,-\gamma_1,0)$, with
    anchor $s$. The only zeros of $\chi_1$ with parameter below $s$ are $\rho_1$ and, for type
    $\mathrm{rc}$, $\overline{\rho_1}$, since every other occurrence has parameter at least $\lambda'\ge p^*\ge s$.
    With $\cS=\{\rho_1\}$ (and $\overline{\rho_1}$ for type $\mathrm{rc}$) nothing remains in the hypotheses of
    \cref{lem:response}. As conjugate zeros of a real character have the same multiplicity $m$,
    \[
      \sum_{\rho\in\cS}m_\rho\Re F(\cdot)=m\bigl(F(\lambda_1-s)+\1_{\mathrm{rc}}\Re F(\lambda_1-s+2i\mu_1)\bigr)
      \ge m\bigl(F(b-s)-\1_{\mathrm{rc}}C_Z\bigr),
    \]
    with $C_Z=\mathcal C_c(s-a)$ from \cref{lem:parabolic}(v) ($0$ if $s\le a$). If $F(b-s)-\1_{\mathrm{rc}}C_Z\ge0$
    this is at least $F(b-s)-\1_{\mathrm{rc}}C_Z$, and we take $r=F(b-s)-\1_{\mathrm{rc}}C_Z-\frac16f(0)$.
    Otherwise the proposed bound $F(b-s)-C_Z-\frac16f(0)$ is negative and the entry receives the feature
    $0$; this is admissible whatever the true response bound $r_j$ is, by the last clause of
    \cref{cor:zeroform} (with $v'_j=0$). The diagonal numerator is $D(0)-d_1$.
  \item \emph{Reserved family.} Entries $(\chi,\gamma_\chi,s-\sigma_2)$ for its $n_2$ characters at their
    height-one representatives, with $\sigma_2=\min(\max(a,\min(l_2,\lo_2)),s)$. Every zero of these
    characters in $R(l)$ has parameter at least $\max(\lambda_1,\lambda_2)\ge\max(a,l_2)\ge\sigma_2$, so there is no
    zero to the right, $r=F(\hi_2-\sigma_2)-\frac16f(0)$, and the diagonal numerator is $D(s-\sigma_2)-d_1$. With
    second-family columns, the column $[\lo_{2,j},\hi_{2,j}]$ uses $\sigma=\min(\max(a,\min(l_2,\lo_{2,j})),s)$,
    $r=F(\hi_{2,j}-\sigma)-\frac16f(0)$ and the diagonal numerator $D(s-\sigma)-d_1$.
  \item \emph{Outside first family} (family row, $s=s_1$, $\Xi=0$). Global entries $(\chi_1,\gamma_1,0)$ with
    $\cS=\{\rho_1\}$ and, for type $\mathrm{complex}$, $(\bchi_1,-\gamma_1,0)$ with $\cS=\{\overline{\rho_1}\}$ (for
    type $\mathrm{rc}$ there is the single global entry $(\chi_1,\gamma_1,0)$); $r=F(b-s_1)-\frac16f(0)$. And a
    hidden entry $(\chi,\gamma_h,0)$ for each first-family character $\chi$ with a zero in $R_P$, at its zero
    $\rho_h$ of least parameter $t_\chi\in[\lo_h,\hi_h)$, with $\cS=\{\rho_h\}$ and $r=F(\hi_h-s_1)-\frac16f(0)$.
    No zero to the right. Entries of the same character have normalized height differences at least
    $C_B-C_P>M$, so their pair excess vanishes, and every diagonal numerator is $D(0)-d_1$.
\end{enumerate}

\subsection{The rows}\label{subsec:rowkinds}
Each leaf program has one or more rows, of the kinds listed in \cref{tab:rowkinds}. In the family
row of an inside leaf the first family is entered by (b), by (c) for type $\mathrm{rc}$ (with the
height split of \cref{subsec:subdivision}), or by (d) in a piece of the second-zero split. In the
shifted row every ordinary anchor is at most $s\le l_2$. In the graded rows all characters are
distinct.

\begin{table}[ht]
\centering
\small
\begin{tabular}{@{}lllll@{}}
\toprule
Row & Reference anchor $s$ & Sieve weights & Entries, inside & Entries, outside\\
\midrule
family & $s_1$ & none ($\Xi=0$) & (a), (f) and (b), (c) or (d) & (a), (f), (g)\\
shifted & $\min(1.9,p^*,l_2)$ & none ($\Xi=0$) & (a), (e), (f) & not used\\
graded & $1.5$, $1.6$, $1.9$, $2.1$ or $2.3$ & present & (a), (b), (f) & (a), (f)\\
two-test & the shift of the leaf & two test functions & \cref{thm:twotest}, on $37$ roots & not used\\
\bottomrule
\end{tabular}
\caption{The kinds of near rows. The shifted row is used only if $s>s_1$ and $s\ge a$.}
\label{tab:rowkinds}
\end{table}

In every family, shifted and graded row the safe-anchor hypothesis holds by (Z1) and
\cref{inp:rectangle}, as explained after \cref{def:safeanchor}, and the characters are pairwise distinct
whenever $\Xi\not\equiv0$. All the
entries have heights at most $l$. The first family is entered once per character in a row with
sieve weights (for type $\mathrm{rc}$ only $\rho_1$ is kept); the pair and second-zero entries occur only
in rows with $\Xi=0$.

\subsection{The two-test row}\label{subsec:twotest}
On $37$ inside roots the graded rows alone do not certify, and the leaf program contains in
addition the following near row. It is not an instance of \cref{thm:graded}: its Gram form is at
the shift $s$ of the leaf, which may exceed $\lambda_1$, and it pays explicit corrections for the
quotient characters $\chi_1,\bchi_1$. Its proof follows Xylouris's proof of his Lemma~5.3
\cite[(5.34)--(5.40), pp.~72--76]{Xylouris2011nullstellen}, with two test functions.

Fix $\gamma_G>\gamma_Z>0$ and let $f=f_{\gamma_Z}$ and $g=f_{\gamma_G}$ as in \eqref{eq:ftest}, so that $g>0$ on the
support of $f$. Let $s\le\min(p^*,l_2)$ be the shift of the leaf, and put
\begin{align*}
  R_G&=G(-s),\quad R_Z=\int_0^{2\gamma_Z}\frac{f(t)^2}{g(t)}e^{st}\,dt,\quad \mathcal N=\sqrt{R_GR_Z},\quad
  d_0=\frac{g(0)}{6R_G},\\
  C_G&=\sup_{\Re z\ge a-s}(-\Re G(z)),\quad C_F=\sup_{\Re z\ge a-s}(-\Re F(z)),\quad c_G=\frac{C_G}{R_G},\quad k_\eta=1+\eta .
\end{align*}
The function $f^2/g$, extended by $0$, satisfies Condition~1, since it vanishes to order $6$ at
$2\gamma_Z$. Define the features
\begin{gather*}
  v(\lambda)=\frac{F(\lambda-s)-\frac16f(0)}{\mathcal N}-\eta,\qquad
  v_2=\frac{F(\hi_2-s)-\frac{\phi_2}2f(0)}{\mathcal N}-\eta,\\
  v_f=\frac{F(b-s)-\frac{\phi_1}2f(0)-\1_{\mathrm{rc}}C_F}{\mathcal N}-\eta,
\end{gather*}
with $\phi_1=\frac14$ for a real $\chi_1$ and $\phi_1=\frac13$ otherwise, $\phi_2=\frac14$ if $n_2=1$ and
$\phi_2=\frac13$ otherwise, and $\1_{\mathrm{rc}}$ the indicator of type $\mathrm{rc}$; the correlation term $d=k_\eta(d_0+\eta)$; and the diagonals
\[
  \begin{array}{lll}
    \text{type} & D_o/k_\eta & D_f/k_\eta\\
    \mathrm{rr} & 1-d_0+c_G & 1-d_0\\
    \mathrm{rc} & 1-d_0+2c_G & 1-d_0\\
    \mathrm{complex} & 1-d_0+2c_G & 1-d_0+c_G
  \end{array}
\]

\begin{theorem}[Two-test inequality]\label{thm:twotest}
Use the two-test data defined above, with $\gamma_G>\gamma_Z>0$ and
$s\le\min(p^*,l_2)$. For every sufficiently large $q$ and every configuration of an
inside leaf, form the following entries:
\begin{enumerate}[(i)]
  \item any finite set of distinct ordinary characters, each at a zero
    $\rho_\chi\in R(l)$ with $|\Im\rho_\chi|\le1$;
  \item the first character at $\rho_1$, and also $\bchi_1$ at $\overline{\rho_1}$ when
    $\chi_1$ is nonreal;
  \item the $n_2$ characters of the reserved family at their height-one representatives,
    if a family is reserved.
\end{enumerate}
For type $\mathrm{rc}$, the single first-family entry retains both $\rho_1$ and
$\overline{\rho_1}$. Assign features $v(\lambda_{\rho_\chi})$, $v_f$, and $v_2$ to
groups (i)--(iii), and diagonals $D_o$, $D_f$, and $D_o$, respectively. Then
for every choice of real coefficients $a_j\ge0$,
\[
  \Bigl(\sum_ja_jv_j\Bigr)_+^2\le\sum_jD_ja_j^2+d\Bigl(\sum_ja_j\Bigr)^2 .
\]
\end{theorem}

\begin{proof}
Put $\sigma_s=1-s/\LL$ and, in the Hilbert space of sequences indexed by $n$,
\[
  A_j(n)=\sqrt{\Lambda(n)n^{-\sigma_s}g(t_n)}\,\chi_j(n)n^{-i\gamma_j},\qquad
  B(n)=\chi_0(n)\sqrt{\Lambda(n)n^{-\sigma_s}}\,\frac{f(t_n)}{\sqrt{g(t_n)}} .
\]
Then $-\Re\langle A_j,B\rangle=\LL R_j$, where $R_j$ is the response of the entry at the anchor $s$, and
$\LL\sum_ja_jR_j\le\|\sum_ja_jA_j\|\,\|B\|$.

\emph{Responses.} We use \cref{lem:response} (or directly \cref{inp:X32}). For an ordinary or reserved
character every zero in the disc has parameter at least $s$ (\cref{lem:noright}(b), as $s\le l_2$), so
$R_j\ge F(\lambda_{\rho_\chi}-s)-\frac{\varphi}2f(0)-\eps$. For the first family, the zeros of $\chi_1$ with
parameter below $s$ are $\rho_1$ and, for type $\mathrm{rc}$, $\overline{\rho_1}$ (the others have parameter at
least $\lambda'\ge p^*\ge s$); keeping them gives $F(\lambda_1-s)\ge F(b-s)$ and, for type $\mathrm{rc}$,
$\Re F(\lambda_1-s+2i\mu_1)\ge-C_F$ (\cref{lem:parabolic}(v)). Conjugate zeros of a real character have the
same multiplicity $m\ge1$, so the kept terms total at least $m(F(b-s)-\1_{\mathrm{rc}}C_F)$, which is at least
$F(b-s)-\1_{\mathrm{rc}}C_F$ when this is nonnegative (otherwise $v_f<0$, see below). Since
$\varphi(\chi)=\frac14$ for real $\chi$ and $\varphi(\chi)\le\frac13$ always (\cref{inp:X32}), the terms $\frac{\varphi}2f(0)$ are at
most $\frac16f(0)$, $\frac{\phi_1}2f(0)$ and $\frac{\phi_2}2f(0)$ respectively. Thus
$R_j\ge \mathcal N(v_j+\eta)-\eps$ for entries with positive feature. Entries with nonpositive
feature may be omitted in estimating the positive part on the left; we justify
restoring their coefficients to the right-hand side below.

\emph{Norms.} By \cref{inp:X31} applied to $f^2/g$, $\|B\|^2=\LL R_Z(1+o(1))$, and $\|A_j\|^2=\LL R_G(1+o(1))$.

\emph{Off-diagonal terms.} For $j\ne k$ the characters are distinct and $\langle A_j,A_k\rangle$ is the
$g$-weighted sum for the nonprincipal character $\chi_{jk}=\chi_j\bchi_k$ at $\sigma_s+i(\gamma_j-\gamma_k)$. By \cref{inp:X32}
its real part is at most $\LL(d_0R_G+\eps)$ plus $\LL$ times the sum of $-\Re G(\cdot)$ over the zeros of $\chi_{jk}$ in
the disc with parameter below $s$. If $\chi_{jk}\notin\{\chi_1,\bchi_1\}$ there are none, by \cref{lem:noright}(b),
which applies because $|\gamma_j-\gamma_k|\le2\le2l$.
If $\chi_{jk}=\chi_1$ (or $\bchi_1$), they are $\rho_1$ (resp.\ $\overline{\rho_1}$), and for type $\mathrm{rc}$ both, each
contributing at most $C_G$; these zeros are simple, since a multiple $\rho_1$ would give $\lambda_1=\lambda'\ge s$. For a given $j$, the $k$ with $\chi_j\bchi_k\in\{\chi_1,\bchi_1\}$ are
$\chi_k\in\{\chi_j\bchi_1,\chi_j\chi_1\}$: at most two, and at most one if $\chi_1$ is real; a first-family
character has none if $\chi_1$ is real ($\chi_1\bchi_1=\chi_0$) and at most one ($\chi_1^2$) otherwise. With
$2a_ja_k\le a_j^2+a_k^2$ for each such pair, these corrections give the diagonals of the table.

\emph{Assembly.} For $a\ge0$,
$\|\sum_ja_jA_j\|^2\le\LL R_G[\sum_ja_j^2(1-d_0+\mathrm{exc}_j)+(d_0+\eps)(\sum_ja_j)^2]$, and
$\|B\|^2\le(1+\eps)\LL R_Z$. Dividing by $\LL^2\mathcal N^2$ and absorbing the errors in $\eta$ gives the claim
for the positive-feature entries. To restore any omitted entries, expand the
right-hand side as
$\sum_j(D_j+d)a_j^2+2d\sum_{j<k}a_ja_k$. Its coefficients are nonnegative:
the table gives $D_j+d\ge k_\eta(1+\eta)>0$ and $d>0$.
Restoring nonnegative $a_j$ therefore increases this bound, while a nonpositive feature
can only decrease the positive part on the left.
\end{proof}

\begin{proposition}[Detector mixture]\label{prop:mixture}
\Cref{thm:twotest} remains true when $f$ is replaced by $f_\epsilon=f+\epsilon g$ with $\epsilon\ge0$, provided
$F$, $f(0)$ and $C_F$ are replaced by $F+\epsilon G$, $f(0)+\epsilon g(0)$ and $C_F+\epsilon C_G$, and $\mathcal N$ by
$\mathcal N_\epsilon=\sqrt{R_G(R_Z+2\epsilon F(-s)+\epsilon^2R_G)}$; the Gram quantities are unchanged.
\end{proposition}

\begin{proof}
$f_\epsilon$ satisfies Conditions~1 and~2, and $f_\epsilon^2/g=f^2/g+2\epsilon f+\epsilon^2g$ is a sum of functions
satisfying Condition~1, so $\|B_\epsilon\|^2=\LL(R_Z+2\epsilon F(-s)+\epsilon^2R_G)(1+o(1))$. The rest of the proof
is unchanged, using $-\Re(F+\epsilon G)\le C_F+\epsilon C_G$ on $\Re z\ge a-s$.
\end{proof}

\begin{proposition}[Paired first family]\label{prop:paired}
Use the data and inside-leaf hypotheses of \cref{thm:twotest}. Suppose $\chi_1$ is
nonreal and the specification $\mathfrak s$ has a finite gap
$\lambda'\in[p^*,h]$, where $h=\hi'$. Fix $\zeta>0$ and
a number $\mathcal B$ with
\begin{equation}\label{eq:pairB}
  \frac{\Re G(-s+i\Delta)}{R_G}-\zeta\,\frac{\Re F(\lambda_1-s+i\Delta)+\Re F(\lambda'-s+i\Delta)}{2\mathcal N}\le\mathcal B
\end{equation}
for all real $\Delta$, $\lambda_1\in[a,b]$ and $\lambda'\in[p^*,h]$. Let $e\ge C_F/(2\mathcal N)$, and put
\[
  m=\frac{F(b-s)+F(h-s)}{2\mathcal N}-\frac{f(0)}{6\mathcal N}-\eta-e,\qquad
  D_*=k_\eta\Bigl\{\frac{1+\mathcal B}2-d_0+c_G-\frac{\zeta e}2\Bigr\},\qquad\zeta_*=\frac{k_\eta\zeta}2 .
\]
Let $m'\le m$ and $D'\ge D_*$ with $D'>0$ and $2D'\ge\zeta_*\max(m',0)$. Then, for every configuration
of the leaf, there is $\tau\in[0,\sqrt d]$ satisfying \eqref{eq:T} for the entries of \cref{thm:twotest},
with the first-family feature $m'$ and diagonal $D'$ in place of $v_f$ and $D_f$.
\end{proposition}

\begin{proof}
For $\chi_1$ use the averaged vector $A_f=\frac12(A_{\chi_1,\gamma_1}+A_{\chi_1,\gamma'})$, with an admissible $\rho'$ of height
$\gamma'$ (\cref{sec:notation}; so $\rho'$ is a zero of $\chi_1$ and $|\gamma'|\le l$), and its conjugate for $\bchi_1$. Its
response is the average of the two responses, which we bound by \eqref{eq:X34} at the points
$1-s/\LL+i\gamma_1$ and $1-s/\LL+i\gamma'$; they lie in $R(9l)$, and \eqref{eq:X34} does not require the kept zeros
to lie in a disc. The set $A_1$ of zeros of $\chi_1$ in $R(l)$ with parameter below $s$ is contained in
$\{\rho_1\}$, since every occurrence in $R(l)$ other than the distinguished one has parameter at least
$\lambda'\ge p^*\ge s$. We put $\rho'$ into $A_2$, and also $\rho_1$ if $\lambda_1\ge s$, with their multiplicities (if
$\rho'=\rho_1$ this is $\rho_1$ with multiplicity at least $2$); by \cref{inp:logfree} their number is bounded
independently of $q$. The multiplicities can only increase the response: the terms of $\rho'$ are
$\ge0$ by Condition~2 because $\lambda'\ge s$, the term $F(\lambda_1-s)$ is positive, and if $\rho_1$ is a multiple
zero then $\lambda_1=\lambda'\ge s$, so its terms are $\ge0$ as well. With $\Delta=\LL(\gamma_1-\gamma')$, $P=\Re G(-s+i\Delta)/R_G$ and $Q$ the second
fraction in \eqref{eq:pairB}, the feature of $A_f$ is therefore at least
$\frac{F(b-s)+F(h-s)}{2\mathcal N}-\frac{f(0)}{6\mathcal N}+Q-\eta$, and $\|A_f\|^2=\LL R_G\frac{1+P}2(1+o(1))$ by \cref{inp:X31}. Each
off-diagonal term $\langle A_f,A_k\rangle$ is the average of two terms of the kind estimated in the proof of
\cref{thm:twotest}, at the heights $\gamma_1-\gamma_k$ and $\gamma'-\gamma_k$, whose absolute values are at most
$l+1\le2l$; the term between $A_f$ and its conjugate is an average of four terms, at heights
$2\gamma_1$, $\gamma_1+\gamma'$ and $2\gamma'$, which are also at most $2l$ in absolute value. So \cref{lem:noright}
applies to all of them, averaging creates no new quotient characters, and the off-diagonal estimates
are unchanged (for a cubic $\chi_1$ the quotient $\chi_1^2=\bchi_1$ of the conjugate pair is paid by the $c_G$
in $D_*$). Put $U=Q+e$; then $U\ge0$, since $\lambda'\ge s$ gives
$Q\ge-C_F/(2\mathcal N)$.

By \eqref{eq:pairB}, $P\le\mathcal B+\zeta Q$, so the actual first-family feature is at least $m+U$ and the
actual diagonal $k_\eta\{\frac{1+P}2-d_0+c_G\}$ is at most $D_*+\zeta_*U$. Put $U'=U+(m-m')\ge U\ge0$. Then the actual
feature is at least $m'+U'$, and the actual diagonal is at most $D'+\zeta_*U'$, which is positive.
Replacing the first-family features and diagonals by these bounds preserves the quadratic
inequality of \cref{thm:twotest}, and \cref{prop:threshold} gives $\tau\in[0,\sqrt d]$ satisfying \eqref{eq:T}
with the first-family data $(m'+U',D'+\zeta_*U')$. Finally, for $\tau\ge0$ and $U'\ge0$,
\[
  \frac{(m'+U'-\tau)_+^2}{D'+\zeta_*U'}\ge\frac{(m'-\tau)_+^2}{D'} .
\]
Indeed the right side is $0$ if $m'\le\tau$. Otherwise put $z=m'-\tau\in(0,m']$. The derivative of
$(z+u)^2/(D'+\zeta_*u)$ in $u\ge0$ is $(z+u)(2D'-\zeta_*z+\zeta_*u)/(D'+\zeta_*u)^2\ge0$, because
$2D'\ge\zeta_*m'\ge\zeta_*z$. So \eqref{eq:T} holds at the same $\tau$ with $(m',D')$.
\end{proof}

In the leaf program the two-test row is the threshold form \eqref{eq:T} of \cref{thm:twotest} (or of one of
its two variants, \cref{prop:mixture,prop:paired}):
\[
  \sum_ix_i\frac{(V_i-\tau)_+^2}{D_o}+n\frac{(v_f-\tau)_+^2}{D_f}+n_2\frac{(v_2-\tau)_+^2}{D_o}\le1-\frac{\tau^2}d,
\]
with $V_i=v(\hi_i)$ the feature of the right end of the ordinary bin $i$ (the tail and the
reserved columns have feature $0$ in this row), and negative features replaced by $0$. The
$T^*$-representatives are admissible entries, since \cref{thm:twotest} allows any zero of height at
most $1$, and all of them have parameter at least $l_2\ge s$ (Z3). The reserved family appears as a
fixed term with the feature $v_2$ at $\hi_2$, which bounds its feature over the whole reserved range,
and with the diagonal $D_o$. The row is used on $37$ roots, once on each: with a single test on $31$
of them, with a mixture on $2$, and in the paired form on $4$ (all with a nonreal $\chi_1$, as
\cref{prop:paired} requires). In the paired form the certified pair $(m',D')$ consists of a lower
bound for $m$ and an upper bound for $D_*$, and the condition $2D'\ge\zeta_*\max(m',0)$ is checked for
this pair; the margins $2D'-\zeta_*m'$ are between $0.18$ and $0.30$. The condition \eqref{eq:pairB} is
verified with interval arithmetic. Both of its terms are even in $\Delta$, so $\Delta\ge0$ suffices. On the grid
$\Delta\in\frac1{200}\ZZ\cap[0,30]$ the left side is enclosed at the four corners of the parameter box
$[a,b]\times[p^*,h]$. As the left side is a sum of a function of $\lambda_1$ and a function of $\lambda'$, its
maximum over the box exceeds the maximum over the corners by at most
$\frac\zeta{2\mathcal N}\bigl(\frac{(b-a)^2}8\int t^2f(t)e^{(s-a)_+t}\,dt+\frac{(h-p^*)^2}8\int t^2f(t)\,dt\bigr)$, the linear
interpolation error in each variable, and this is added. Between grid points a
second-derivative bound in $\Delta$ is used, and for $\Delta\ge30$ a closed-form bound from
\cref{lem:parabolic}(iv). The value $\mathcal B$ used is at least $0$.

\subsection{Validity of the rows}

\begin{proposition}[Every configuration satisfies the near rows]\label{prop:rowsvalid}
Let $\mathfrak s$ specify a leaf or one of its subcases. For sufficiently large $q$,
take any configuration in $C(\mathfrak s)$ and assign its characters to columns as in
\cref{prop:realization}, obtaining column values $x$. Then, for every near row $k$,
there is a threshold $\tau_k$ such that
\[
  0\le\tau_k\le\sqrt{d_k/S},\qquad
  \Phi_k(x,\tau_k)\le1-\frac{\tau_k^2}{d_k/S}.
\]
Here $d_k$ is the stored correlation term, $S=10^{16}$, and $\Phi_k$ is the row expression of
\cref{def:leafprogram}. Each row may have its own threshold.
\end{proposition}

\begin{proof}
For a graded, family or shifted row, take as entries the characters of the configuration that
are placed in the columns and family terms of the row: the ordinary characters of $O$ at their
$T^*$-representatives, the reserved characters, and the first family, with the entries of
\cref{subsec:entries}. These form an admissible family of entries for \cref{thm:graded}, with the
kept sets verified in \cref{subsec:entries}, so \cref{cor:zeroform} gives a threshold $\tau$ at
which \eqref{eq:T} holds with the normalized features, diagonals and correlation term. Each column's feature is
at most the feature of every character placed in it (features are taken at the right end of the
column, and $F$ is decreasing), or is $0$; each column's diagonal is at least theirs, or the column
has feature $0$ (an entry whose diagonal numerator is not safely positive is omitted, and its column
is given feature $0$), in which case it contributes nothing for $\tau\ge0$; each family
term $(n,v,D)$ receives $n$ entries with feature at least $v$ and diagonal at most $D$; and the
characters in the tail column and in the hidden tail are simply omitted from the row. Since each term $(v-\tau)_+^2/D$
increases with $v$ and decreases with $D$ for $\tau\ge0$, the row of the leaf holds at $\tau$. For the
two-test row the same argument applies to \cref{thm:twotest} and its variants, through
\cref{prop:threshold}.
\end{proof}

The numbers $I_u$, $D_u$, $D(\delta)$, the transforms and the special terms $R_{\lo}$, $c^{\hi}$, $R_p$, $R_q$,
$C_Z$ and $C_F$ are computed from the rational parameters of each row with outward interval arithmetic. The
transforms are enclosed with the closed form of \cref{lem:parabolic}(iv) or with its power series near
$0$. The infima and suprema over heights in $R_{\lo}$, $c^{\hi}$, $R_p$ and $R_q$ are enclosed on grids of step
at most $\frac1{50}$, with second-derivative bounds $\int t^2\hat f_c(t)e^{dt}\,dt$ between grid points and
Lipschitz bounds in the real parts, and by the tail bounds of \cref{lem:parabolic}(iv) beyond the
grid. The constants $C_Z$, $C_F$ and $C_G$ are enclosed on the line $\Re z=-d$ (by
\cref{lem:parabolic}(v)), for heights in $[0,20]$ by adaptive linear interpolation, starting from
$200$ cells of width $\frac1{10}$ and bisecting cells until the largest cell bound is within a small
fixed tolerance of the best value found, with the same second-derivative bound, and beyond height
$20$ by a closed-form tail bound. The features, diagonals and correlation terms of all rows other than the $37$ two-test rows have been
recomputed from their rational parameters by an independent program whose soundness is formally
verified (\cref{sec:verification}).
